% AVOI -- manuscript-ready results text (conservative). Numbers from results/stage2/tables and results/stage1_ungated.
% Stage 1 = original ungated experiment (historical). Stage 2 = evidence-model adequacy safety condition + confirmation on new data.
\section{Results}

\subsection{Experimental history and design}
The evaluation proceeded in two stages. Stage~1 was the original, pre-registered experiment (hypotheses H1--H5). It had
no evidence-model adequacy condition, and its results are reported unchanged. In Stage~1 the evidence model transferred
poorly between AIT-LDS testbeds. This motivated an evidence-model adequacy safety condition, which was defined for
Stage~2 and was not part of the Stage-1 pre-registration. Stage~2 was carried out in a fixed order:
\begin{enumerate}
  \item The adequacy policy, the Stage-2 hypotheses (S1--S4), the data partitions and all thresholds were frozen.
  \item Evidence models and calibrators were selected on development data only (the Stage-1 data).
  \item Adequacy was assessed on a separate validation partition of new data, and the decisions were frozen.
  \item The methods were evaluated on a disjoint confirmation partition.
\end{enumerate}
The new data were live IAM scenarios generated with a new seed, AIT-LDS testbeds not used in Stage~1, and CERT release
r5.2. The AVOI decision rule is unchanged from Stage~1. Unsafe assurance rate (UAR) is
$\#(\text{Assured}\wedge Y{=}0)/\#(Y{=}0)$, and every UAR is reported together with autonomous coverage.

\subsection{Stage 1}
AVOI substantially reduced the number of evidence acquisitions relative to exhaustive and fixed-order strategies on all
three datasets. It did not uniformly reduce unsafe assurance relative to information gain (H2 not supported). H1 was not
supported as registered, because its matched-safety definition was degenerate. Exhaustive evidence collection achieved
lower expected loss on the controlled IAM testbed (0.316 vs.\ 0.381). On AIT-LDS, all methods showed high unsafe
assurance (UAR $\approx 0.31$).

\subsection{Evidence-model adequacy}
Four of nine evidence models satisfied the predefined autonomous-assurance eligibility policy, all four on the controlled
IAM testbed. The remaining five failed as follows:
\begin{itemize}
  \item AIT exfiltration model: its discrimination was inverted on the new testbeds (AUROC 0.024, 95\% CI 0.013--0.038),
    so it was classified UNSUPPORTED.
  \item AIT intrusion model: it lacked sample support (nine validation violations).
  \item CERT r5.2 model: it was over-confident (calibration slope 0.28, ECE 0.18).
  \item Two IAM models: they discriminated almost perfectly but had calibration slopes above the predefined upper bound
    of 1.5.
\end{itemize}
Evidence-model adequacy therefore emerged as an important boundary condition for autonomous value-guided acquisition.
These criteria are the eligibility rules of this study, not universal standards, and passing them does not demonstrate
regulatory compliance.

\subsection{Acquisition policy on validated evidence models}
On confirmation episodes of validated models (768 IAM episodes), AVOI required 1.45 acquisitions per episode. This was
4.55 fewer than exhaustive collection and 3.16 fewer than a fixed checklist (paired Wilcoxon, Holm-adjusted
$p<10^{-128}$). It was also 1.08 fewer than information gain. Unsafe assurance was 0.003 for AVOI, identical to
information gain and lookahead-2. Expected loss was higher than with exhaustive evidence (difference $+0.042$, 95\% CI
0.005--0.082) and numerically, but not significantly, higher than with information gain ($+0.017$, $-0.001$--0.035).
AVOI's normalised acquisition cost was lower than that of exhaustive, fixed-order and information-gain strategies, but
higher than that of the lookahead-2 policy. Within each IAM objective, AVOI acquired sources in the same order in every
episode; the savings arose from value-based stopping rather than from adaptive choice of source.

The asymmetric assurance-loss formulation reduced unsafe assurance relative to symmetric generic VOI (0.003 vs.\ 0.268;
exact McNemar, Holm-adjusted $p=2.4\times10^{-22}$). Part of this difference reflects AVOI's escalation option. Against
an otherwise identical symmetric-loss variant, it lowered UAR by 0.021 (95\% CI 0.004--0.039), while the paired McNemar
test on discordant unsafe assurances was not significant ($p=0.07$). Along the risk--coverage curve, AVOI's trade-off was
markedly better than confidence-threshold stopping (normalised area 0.0023 vs.\ 0.097) and similar to information gain.

\subsection{Safety gating}
Safety gating reduced autonomy where the underlying evidence model failed the predefined adequacy policy. The gate was
applied identically to all methods. It produced no autonomous decisions for any non-validated model, and it never
converted model inadequacy into a negative assurance decision. Every change in unsafe assurance was accompanied by a loss
of autonomous coverage:
\begin{itemize}
  \item Across the IAM testbed, UAR fell from 0.013 to 0.002, while coverage fell from 0.785 to 0.451.
  \item For the two IAM models restricted by the calibration-slope rule, the gate prevented five false assurances but
    withheld 346 correct autonomous decisions.
  \item For the AIT exfiltration and CERT models, the ungated policy had already escalated every confirmation episode
    at its validation-tuned threshold, so the gate changed no decision.
\end{itemize}
These reductions reflect withheld autonomy, not an improved acquisition policy.

\subsection{Limitations of the evidence}
The prospectively defined joint criterion (UAR $\le\eta$ and coverage $\ge\gamma$) remained degenerate: policies that
issue only negative decisions without evidence satisfied it. We therefore make no claim of efficiency at matched safety.
Evidence-uncertainty escalation on validated models prevented false non-assurances but no false assurances. All eligible
models come from a controlled testbed, and the confirmation data share generators or benchmark families with the
development data. AVOI should therefore not currently be used autonomously on AIT-LDS or CERT-style evidence with these
evidence models, and it is not a replacement for human assurance review.
